\documentclass[../paper.tex]{subfiles}

\begin{document}

We use the term algorithm to denote a concrete computational procedure. Except where explicitly stated, the algorithms considered here are heuristics for modularity maximization and do not guarantee optimality.

Let $G = (V, E)$ be an undirected simple graph, with vertex set $V$ and edge set $E$. We use $n := |V|$ and $m := |E|$ to denote the number of vertices and edges. The \emph{neighborhood} of $u$ is $N(u) := \{v : \{u, v\} \in E\}$ and the \emph{degree} is $d(u) := |N(u)|$. A \textit{clustering} $C$ of a graph is the partition of the nodes $V$ into $k$ disjoint blocks $V_1,...,V_k$ such that the union $V_1 \cup ... \cup V_k = V$.

The \emph{modularity} score of a clustering is the fraction of edges that are within one cluster minus the expected fraction if the edges were distributed at random. Let $L_i := |\{\{u, v\} \in E : \{u, v\} \subseteq V_i \}|$ be the set of edges within block $i$. With that, we can define modularity as:

\[Q(C) := \sum_{i = 1}^{k}{\left(\frac{L_i}{m} - \left(\sum_{u \in V_i}{\frac{d(u)}{2m}}\right)^2 \right)} \]

\ifSubfilesClassLoaded{%
    \bibliography{../paper.bib}%
}{}

\end{document}